% lem-separated-brackets: limit-shapes.tex:1532-1539 (label lem:separated-brackets)
\paragraph{Paper excerpt.}
\begin{verbatim}
\begin{lemma}\label{lem:separated-brackets}
Let $\gauge$ be the Euclidean norm, and let $0<\drift<\nf1d$. There exist $0<c_0(d,\drift)\leq C_0(d,\drift)$ and $C(d,\drift)<\infty$ such that, for all sufficiently large~$n$, on~$\Omega_n$, for $1\leq i\leq m$ and $1\leq j\leq m$ with $j\neq i$,
\begin{equation*}
c_0\leq\langle S^i\rangle_n\leq C_0
\qquad\text{and}\qquad
|\langle S^i,S^j\rangle_n|\leq Cr_n^{-\nf14}\, .
\end{equation*}
\end{lemma}
\end{verbatim}
\paragraph{Lean statement.}
\begin{verbatim}
theorem CERW.Frozen.separated_brackets {d : ℕ} (hd : 2 ≤ d) :
    let ωd : ℝ := (volume (Metric.ball (0 : EuclideanSpace ℝ (Fin d)) 1)).toReal
    ∀ ε : ℝ, 0 < ε → ε < 1 / (d : ℝ) →
    let r : ℕ → ℝ := fun n => ((d + 1) * n / (2 * d * ε * ωd)) ^ ((1 : ℝ) / (d + 1))
    let e₁ : Site d := LatticeProb.unit ⟨0, by omega⟩
    let g : Site d → ℝ := CERW.latticeKernel d
    let σ : ℕ → ℝ := fun n =>
      if d = 2 then Real.sqrt (r n * Real.log (r n)) else Real.sqrt (r n)
    let k : ℕ → ℕ := fun n => ⌊r n ^ ((1 : ℝ) / 8)⌋₊
    let m : ℕ → ℕ := fun n => ⌊r n ^ ((1 : ℝ) / 8)⌋₊
    let y : ℕ → ℕ → Site d := fun n i => ((i * ⌈r n ^ ((1 : ℝ) / 4)⌉₊ : ℕ) : ℤ) • e₁
    let f : ℕ → ℕ → Site d → ℝ := fun n i z =>
      if d = 2 then g (z - (y n i + (k n : ℤ) • e₁)) - g (z - y n i) else g (z - y n i)
    ∃ c₀ C₀ C : ℝ, 0 < c₀ ∧ c₀ ≤ C₀ ∧ 0 < C ∧ ∃ n₀ : ℕ,
      ∀ {Ω : Type u} [MeasurableSpace Ω] (μ : Measure Ω) [IsProbabilityMeasure μ]
        (X : ℕ → Ω → Site d), CERW.IsCERW μ ε X → ∀ n : ℕ, n₀ ≤ n →
        μ {ω | ¬ ∀ i j : ℕ, 1 ≤ i → i ≤ m n → 1 ≤ j → j ≤ m n →
            let B : ℝ := CERW.dynkinBracket (CERW.stepProb d ε) (f n i) (f n j) (X · ω) n /
              σ n ^ 2
            (i = j → c₀ ≤ B ∧ B ≤ C₀) ∧ (i ≠ j → |B| ≤ C * r n ^ (-(1 : ℝ) / 4))}
          ≤ ENNReal.ofReal (C * (n : ℝ) ^ (-(10 : ℝ)))
--
\end{verbatim}
\paragraph{Transcription.}
The statement assumes d ≥ 2 and 0 < ε < 1/d. Constants c₀, C₀, C and the threshold n₀ precede the probability space, with 0 < c₀ ≤ C₀. For n ≥ n₀ the complement of the simultaneous bracket bounds has outer measure at most C n⁻¹⁰. The paper states these bounds on the event Ω_n of Section 5.1. The deterministic implication on that event is proved in CERW/Support/Lower/SeparatedBracketsEvent.lean. The registered probability estimate follows from this implication and the probability bound for Ω_n. The family size, normalization and every pair of indices in the bracket bounds are those in the paper.
